\section{Half-angle, half-side, half-tangent, and haversine formulas}
\label{sec:half-formulas}

Recall
\[
 p=\frac{a+b+c}{2},\qquad P=\frac{A+B+C}{2}.
\]
Every side and angle half lies in \((0,\pi/2)\).  Consequently all square roots
in this section mean the nonnegative root and no independent sign choice is
hidden in the notation.

\subsection{Half-angle formulas}

\begin{theorem}[Spherical half-angle formulas]
\label{thm:half-angle}
For the angle \(A\),
\begin{align}
 \sin^2\frac A2
 &=\frac{\sin(p-b)\sin(p-c)}{\sin b\sin c},                \tag{6.1}\\
 \cos^2\frac A2
 &=\frac{\sin p\sin(p-a)}{\sin b\sin c},                  \tag{6.2}\\
 \tan^2\frac A2
 &=\frac{\sin(p-b)\sin(p-c)}{\sin p\sin(p-a)}.            \tag{6.3}
\end{align}
Hence
\begin{align}
 \sin\frac A2
 &=\sqrt{\frac{\sin(p-b)\sin(p-c)}{\sin b\sin c}},        \tag{6.4}\\
 \cos\frac A2
 &=\sqrt{\frac{\sin p\sin(p-a)}{\sin b\sin c}},           \tag{6.5}\\
 \tan\frac A2
 &=\sqrt{\frac{\sin(p-b)\sin(p-c)}{\sin p\sin(p-a)}}.     \tag{6.6}
\end{align}
The formulas for \(B,C\) are cyclic.
\end{theorem}

\begin{proof}
Subtract the side cosine law from the cosine difference formula:
\begin{align*}
 \cos(b-c)-\cos a
 &=\sin b\sin c(1-\cos A)\\
 &=2\sin b\sin c\sin^2(A/2).
\end{align*}
On the other hand, product-to-sum gives
\[
 \cos(b-c)-\cos a=2\sin(p-b)\sin(p-c).
\]
This proves (6.1).  Similarly,
\begin{align*}
 \cos a-\cos(b+c)
 &=\sin b\sin c(1+\cos A)\\
 &=2\sin b\sin c\cos^2(A/2),
\end{align*}
whereas
\[
 \cos a-\cos(b+c)=2\sin p\sin(p-a).
\]
This proves (6.2).  Divide (6.1) by (6.2), using the positive apartness of
\(\cos(A/2)\), to obtain (6.3).  Positive-root uniqueness gives
(6.4)--(6.6).  The proof uses only the cosine law and finite trigonometric
identities already available in \(\ipra\).
\end{proof}

A useful common-factor form is obtained by setting
\begin{equation}
 k=\frac{1}{\sqrt{\sin p\sin(p-a)\sin(p-b)\sin(p-c)}}.     \tag{6.7}
\end{equation}
Then
\begin{equation}
 \tan\frac A2=k\sin(p-b)\sin(p-c),                         \tag{6.8}
\end{equation}
with cyclic companions.

\subsection{Half-side formulas}

\begin{theorem}[Spherical half-side formulas]
\label{thm:half-side}
For the side \(a\),
\begin{align}
 \sin^2\frac a2
 &=-\frac{\cos P\cos(P-A)}{\sin B\sin C},                 \tag{6.9}\\
 \cos^2\frac a2
 &=\frac{\cos(P-B)\cos(P-C)}{\sin B\sin C},                \tag{6.10}\\
 \tan^2\frac a2
 &=-\frac{\cos P\cos(P-A)}{\cos(P-B)\cos(P-C)}.           \tag{6.11}
\end{align}
Consequently
\begin{equation}
 \tan\frac a2=K\cos(P-A),                                 \tag{6.12}
\end{equation}
where
\begin{equation}
 K=\sqrt{-\frac{\cos P}
 {\cos(P-A)\cos(P-B)\cos(P-C)}}.                          \tag{6.13}
\end{equation}
The formulas for \(b,c\) are cyclic.
\end{theorem}

\begin{proof}
Apply Theorem~\ref{thm:half-angle} to the polar triangle.  Its angle at
\(A^*\) is \(\pi-a\), its opposite side is \(\pi-A\), and its half-sums are
given by (5.7).  Thus
\[
 \cos^2\frac{A^*}{2}=\sin^2\frac a2
 =\frac{\sin p^*\sin(p^*-a^*)}{\sin b^*\sin c^*}
 =-\frac{\cos P\cos(P-A)}{\sin B\sin C},
\]
which is (6.9).  Likewise
\[
 \sin^2\frac{A^*}{2}=\cos^2\frac a2
 =\frac{\cos(P-B)\cos(P-C)}{\sin B\sin C},
\]
giving (6.10).  Their quotient gives (6.11).  Equation (6.12) is simply the
positive square root of (6.11) after extracting the common factor (6.13).
Every sign is certified by the proper-triangle hypotheses and the polar
construction.
\end{proof}

\begin{remark}[The Wikipedia half-side formula]
Equation (6.11), followed by the positive square root, is the standard
half-side formula
\[
 \tan\frac a2=
 \sqrt{-\frac{\cos P\cos(P-A)}
 {\cos(P-B)\cos(P-C)}}.
\]
The constructive statement includes the proof that the radicand is
nonnegative and a branch certificate selecting \(0<a/2<\pi/2\).
\end{remark}

\subsection{Direct vector half-angle formulas}

The vector certificate yields particularly short expressions:
\begin{align}
 \tan\frac A2
 &=\frac{\sin A}{1+\cos A}
 =\frac{\tau}{s_bs_c+x-yz},                                \tag{6.14}\\
 \tan\frac a2
 &=\frac{\sin a}{1+\cos a}
 =\frac{s_a}{1+x}.                                         \tag{6.15}
\end{align}
The denominators are positive because \(A,a\in(0,\pi)\).  These formulas are
well suited to computation: they use one square root for each side sine and no
inverse trigonometric function.

\subsection{Half-tangent rational forms}

Put
\[
 t_a=\tan\frac a2,
 \quad t_b=\tan\frac b2,
 \quad t_c=\tan\frac c2,
 \qquad
 t_A=\tan\frac A2,
 \quad t_B=\tan\frac B2,
 \quad t_C=\tan\frac C2.
\]
Tangent addition, together with Theorems~\ref{thm:half-angle} and
\ref{thm:half-side}, gives the entirely rational forms
\begin{equation}
 t_A^2=
 \frac{(t_a-t_b+t_c+t_at_bt_c)(t_a+t_b-t_c+t_at_bt_c)}
 {(t_a+t_b+t_c-t_at_bt_c)(-t_a+t_b+t_c+t_at_bt_c)},        \tag{6.16}
\end{equation}
and
\begin{equation}
 t_a^2=
 \frac{(t_Bt_C+t_Ct_A+t_At_B-1)
       (-t_Bt_C+t_Ct_A+t_At_B+1)}
 {(t_Bt_C-t_Ct_A+t_At_B+1)
       (t_Bt_C+t_Ct_A-t_At_B+1)}.                          \tag{6.17}
\end{equation}
All displayed denominator factors carry apartness witnesses.  The remaining
five equations are obtained cyclically.  To verify (6.16), write every
\(\sin(p-a)\) and \(\sin p\) using the tangent-of-a-sum formula and cancel the
common positive products \(\prod(1+t_x^2)^{-1/2}\).  Formula (6.17) is its polar
dual.  This derivation is finite rational algebra, so after cross multiplication
it is directly suitable for an exact certificate checker.

\subsection{Haversine formulas}

Define
\begin{equation}
 \hav t=\sin^2\frac t2=\frac{1-\cos t}{2}.                 \tag{6.18}
\end{equation}
The side cosine law becomes the numerically stable navigation formula
\begin{equation}
 \hav a=\hav(b-c)+\sin b\sin c\,\hav A,                    \tag{6.19}
\end{equation}
with cyclic companions.  Indeed,
\[
 \cos(b-c)-\cos a=\sin b\sin c(1-\cos A),
\]
and division by \(2\) gives (6.19).  The dual angle form is
\begin{equation}
 \hav A=\hav(\pi-B-C)+\sin B\sin C\,\hav a.                \tag{6.20}
\end{equation}
When \(\sin b\sin c\apart0\), one may solve for the angle:
\begin{equation}
 \hav A=\frac{\hav a-\hav(b-c)}{\sin b\sin c}.             \tag{6.21}
\end{equation}
Equations (6.18)--(6.21), their cyclic forms, and their polar duals account for
the standard haversine family without any loss of information from the cosine
law.
